\section{Introduction}
\vspace{-0.5em}
\begin{figure}[t]
    \centering
    \includegraphics[width=\linewidth]{figures/teaser.pdf}
    \caption{\textbf{Left}: Zero-shot shape classification on the Objaverse-LVIS (1,156 categories) and ModelNet40 datasets. OpenShape outperforms previous methods by a large margin. We exclude shapes in Objaverse-LVIS during training, and we also retrain ULIP~\cite{xue2022ulip} on our ensembled training shapes for fair comparison. \textbf{Right}: Our shape representations encode a broad range of semantic and visual concepts. We input two 3D shapes and use their shape embeddings to retrieve the top three shapes whose embeddings are simultaneously closest to both inputs. See Section.~\ref{sec:app} for more details.}
    \label{fig:teaser}
    \vspace{-0.7em}
\end{figure}

3D shape understanding has recently garnered a surge of interest driven by the growing demands in real-world applications, such as augmented/virtual reality, autonomous driving, and robotics. Despite significant advancements in 3D recognition and analysis, existing data-driven approaches are still greatly limited by the scale of 3D training datasets and tend to exhibit poor generalization when facing unseen shape categories, hindering the deployment of existing models in real-world applications.

Note that 3D shapes and 2D images can be easily linked through rendering, and the dataset scale issue of 2D images has been remarkably addressed, as shown in recent works such as CLIP~\cite{radford2021learning}. Therefore, many recent studies aim to utilize pre-trained 2D image-language models~\cite{radford2021learning,rombach2022high} to assist 3D tasks, such as 3D generation~\cite{hong2022avatarclip,jain2022zero,michel2022text2mesh,sanghi2022clip,lee2022understanding,canfes2023text} and 3D scene-level segmentation~\cite{ha2022semantic,jatavallabhula2023conceptfusion,ding2022language,yang2023regionplc,lu2023open,peng2022openscene}. Regarding 3D shape-level understanding, a straightforward idea is to project 3D data to the 2D domain through rendering and use CLIP to analyze the 2D images, thereby enabling zero-shot 3D shape classification~\cite{zhang2022pointclip,zhu2022pointclipv2}. However, these methods suffer from occlusion and information loss during projection, and unnecessary latency due to point cloud rendering and multiple CLIP inferences.

To overcome the limitations caused by projection, it is necessary to train a 3D-native model by distilling knowledge from pretrained 2D models. However, training a 3D-native model requires a set of 3D shapes, and the amount of knowledge that can be distilled is determined by the size of the 3D dataset. For example, ULIP~\cite{xue2022ulip} aims to learn a joint representation space between language, 2D images, and 3D shapes, but uses a small-scale 3D dataset ShapeNetCore~\cite{chang2015shapenet} for knowledge distillation. Specifically, ULIP fixes the 2D CLIP text and image encoders and trains a dedicated 3D-native point cloud encoder to extract 3D shape representations. The 3D encoder strives to align the 3D shape embedding space with the CLIP image and language embedding spaces by utilizing contrastive learning across all three modalities. However, since ULIP is only trained on 52K shapes of 55 object categories, it still struggles with out-of-distribution shape categories and fails to demonstrate an impressive open-world understanding of 3D shapes. 

In this work, we propose a novel method called OpenShape, which follows a similar paradigm as ULIP but aims to achieve a more generalized and scalable joint representation space encompassing language, 2D images, and 3D shapes. Our focus mainly lies on scaling up representation learning and addressing corresponding challenges. In OpenShape, we emphasize four key factors during the training process: (a) data scale: we significantly increase the scale of 3D training data by combining four public 3D shape datasets, resulting in 876k 3D shapes covering much more diverse categories; (b) text quality: the 3D shapes from our main dataset, Objaverse~\cite{objaverse}, is dominated with inaccurate or uninformative text descriptions. Given the data scale, we propose three strategies to automatically filter and enrich the text descriptions; (c) 3D backbone scaling: since most existing 3D backbones target small datasets, we find that it's important but non-trivial to scale up the 3D backbones; and (d) data resampling: since the ensembled dataset is highly unbalanced, we utilize hard negative mining to improve the model's discriminative ability.

We first evaluate OpenShape on the zero-shot 3D shape classification task. As shown in Figure~\ref{fig:teaser}, OpenShape outperforms previous zero-shot approaches on the ModelNet40 dataset by at least 20\%. Moreover, OpenShape excels at handling long-tail categories. On the challenging Objaverse-LVIS dataset, which contains 1,156 categories, OpenShape achieves a 46.8\% accuracy, significantly surpassing previous methods. Notably, this performance gap remains even when ULIP is retrained on our ensembled datasets, highlighting the superiority of our text enrichment and training strategies. Besides zero-shot classification, we present demos that showcase the wide range of visual and semantic concepts learned by OpenShape. For example, in Figure~\ref{fig:teaser}-right, we take two 3D shapes as input and use their OpenShape embeddings to retrieve the top three shapes whose embeddings are simultaneously closest to both inputs from our ensembled dataset. The retrieved shapes exhibit an interesting combination of the semantic and geometric elements from both input shapes. Furthermore, since we align our 3D shape embedding space with the CLIP language and image embedding space, we demonstrate that OpenShape embeddings can be easily integrated with other CLIP-based models to perform cross-modality tasks such as point cloud captioning and point cloud-conditioned image generation.
